\subsubsection{Congruence ternaire de la trace et optimalité de l'exposant}
\label{mg09:congruencia}

L'identité de niveau trois permet de comparer tous les coefficients
de la trace construite à ceux de sa réalisation modulaire. Le domaine
approprié pour cette conséquence est l'anneau des séries formelles : chaque
coefficient dépend d'un nombre fini de facteurs du produit,
tandis que l'énoncé conserve tous les degrés.

\begin{lemma}[Unité entière du produit d'ordre trois]
La série
\begin{equation}
 u(q)=q\,t_3(q)=\prod_{n\geq1}(1+q^n+q^{2n})^{-12}
 \label{mg09:unidad}
\end{equation}
appartient à \(1+q\mathbb Z[[q]]\), et son inverse appartient au
même ensemble.
\end{lemma}

\begin{proof}
Chaque polynôme \(1+q^n+q^{2n}\) a pour terme constant un.
Pour toute série \(v(q)=1+\sum_{r\geq1}v_rq^r\) à
coefficients entiers, l'équation \(v(q)w(q)=1\) détermine
récursivement
\[
 w_0=1,\qquad w_r=-\sum_{i=1}^{r}v_iw_{r-i}\quad(r\geq1).
\]
Cette récurrence prouve l'existence, l'unicité et le caractère entier de l'inverse. Appliquée à chaque facteur, elle prouve également le caractère entier de sa puissance inverse douze. En outre, \((1+q^n+q^{2n})^{-12}\in1+q^n\mathbb Z[[q]]\). Pour calculer un coefficient de degré au plus \(r\), les facteurs avec \(n\leq r\) suffisent ; les produits partiels se stabilisent coefficient par coefficient. Ainsi, \eqref{mg09:unidad} définit une série à coefficients entiers et à terme constant un. La même récurrence, appliquée maintenant à \(u\), donne \(u^{-1}\in1+q\mathbb Z[[q]]\).
\end{proof}

\begin{proposition}[Divisibilité uniforme exacte]
Pour la trace \(t_3\) et la série \(J=j-744\) reliées par
\eqref{km:j-nivel3}, on a
\begin{equation}
 J-(t_3+12)\in3^9q\mathbb Z[[q]].
 \label{mg09:congruencia-exacta}
\end{equation}
L'exposant uniforme de divisibilité par trois est exactement neuf :
\begin{equation}
 [q]\bigl(J-t_3-12\bigr)=10\cdot3^9=196830.
 \label{mg09:optimalidad}
\end{equation}
\end{proposition}

\begin{proof}
Comme \(t_3=q^{-1}u\), le lemme donne
\(t_3^{-r}=q^ru^{-r}\in q^r\mathbb Z[[q]]\) pour
\(r=1,2,3\). L'identité \eqref{km:j-nivel3} se transforme en
\begin{equation}
 J-t_3-12
 =3^9q\left(10u^{-1}+4\cdot3^5q\,u^{-2}
                         +3^9q^2u^{-3}\right).
 \label{mg09:factorizacion}
\end{equation}
Tous les termes entre parenthèses sont des séries entières ; cela prouve
\eqref{mg09:congruencia-exacta}. Son terme constant est \(10\),
car \(u(0)=u^{-1}(0)=1\), et les deux derniers termes contiennent
\(q\) et \(q^2\). On obtient \eqref{mg09:optimalidad}. Comme
\(3\nmid10\), le coefficient de \(q\) a une valuation
\(3\)-adique de neuf, ce qui établit l'optimalité de l'exposant uniforme.
\end{proof}

En particulier, \([q^r]J\equiv[q^r]t_3\pmod{3^9}\) pour tout
\(r\geq1\). Le quantificateur procède de l'identité formelle et de
l'inversion entière démontrée dans le lemme. Une vérification jusqu'à
un degré fini quelconque est une spécialisation de cette relation.
La trace procède toujours de l'automorphisme d'ordre trois du réseau
construit à partir de l'incidence HMT ; l'uniformisation agit ensuite
sur cette trace. La congruence conserve cet ordre et détermine une
relation exacte entre ses deux lectures graduées.
